\chapter{Résultats des simulations}
\section{Trajectoires aux prix historiques}
La figure \ref{fig:historique} présente les 18 profils dans les deux périmètres. Les trajectoires utilisent les indices de prix observés du HCP et les coefficients fixes de l'ENNVM. Les variations diffèrent selon le périmètre et la composition des budgets. Entre septembre 2022 et octobre 2026, les hausses de coût sont comprises entre 8,97 et 9,45\,\% pour le périmètre prudent, et entre 7,60 et 8,34\,\% pour l'élargi. Pour l'ensemble national, les hausses sont de 9,22 et 7,85\,\%. Les contributions correspondantes au budget initial sont respectivement de 3,82 et 7,31 points. Le sous-panier prudent augmente davantage en proportion de son propre coût initial. Le périmètre élargi pèse davantage dans le budget total, ce qui conduit à une contribution supérieure.

Les variations d'octobre 2026 sur un mois ou sur douze mois sont disponibles dans le CSV historique. Le tableau \ref{tab:historique} résume les 18 profils à cette date par rapport à la référence. Les trajectoires antérieures à septembre 2022 appliquent rétrospectivement les structures budgétaires de l'enquête.
\figureprojet[1.10\linewidth]{04_cout_panier_historique.png}{Indices de coût par milieu et périmètre, de janvier 2017 à octobre 2026 (septembre 2022 = 100). C1 à C5 désignent les quintiles de dépenses de chaque milieu. Certaines courbes se superposent. Source : notebook 04, d'après le HCP \cite{hcpIndicateurs,hcpENNVM}.}{fig:historique}
\clearpage
\begin{table}[!htbp]
\centering\small
\renewcommand{\arraystretch}{1.10}
\caption{Variation du coût et contribution au budget initial en octobre 2026}\label{tab:historique}
\begin{tabular}{llrrrr}
\hline
Milieu & Classe & Prudent (\%) & Prudent (pts) & Élargi (\%) & Élargi (pts) \\
\hline
National & C1 & 9{,}42 & 4{,}89 & 8{,}21 & 8{,}05 \\
National & C2 & 9{,}38 & 4{,}58 & 8{,}16 & 7{,}92 \\
National & C3 & 9{,}36 & 4{,}32 & 8{,}05 & 7{,}72 \\
National & C4 & 9{,}31 & 4{,}00 & 7{,}90 & 7{,}46 \\
National & C5 & 9{,}04 & 3{,}26 & 7{,}65 & 6{,}89 \\
National & \textbf{Ensemble} & 9{,}22 & 3{,}82 & 7{,}85 & 7{,}31 \\
\noalign{\vskip 3pt}
Urbain & C1 & 9{,}39 & 4{,}23 & 8{,}07 & 7{,}89 \\
Urbain & C2 & 9{,}35 & 4{,}07 & 7{,}99 & 7{,}69 \\
Urbain & C3 & 9{,}31 & 3{,}94 & 7{,}93 & 7{,}54 \\
Urbain & C4 & 9{,}26 & 3{,}74 & 7{,}81 & 7{,}29 \\
Urbain & C5 & 8{,}97 & 3{,}10 & 7{,}60 & 6{,}80 \\
Urbain & \textbf{Ensemble} & 9{,}17 & 3{,}54 & 7{,}77 & 7{,}18 \\
\noalign{\vskip 3pt}
Rural & C1 & 9{,}45 & 5{,}28 & 8{,}34 & 8{,}20 \\
Rural & C2 & 9{,}41 & 5{,}18 & 8{,}26 & 8{,}08 \\
Rural & C3 & 9{,}41 & 5{,}08 & 8{,}29 & 8{,}04 \\
Rural & C4 & 9{,}39 & 4{,}94 & 8{,}22 & 7{,}88 \\
Rural & C5 & 9{,}34 & 4{,}34 & 8{,}00 & 7{,}41 \\
Rural & \textbf{Ensemble} & 9{,}38 & 4{,}78 & 8{,}16 & 7{,}77 \\
\hline
\end{tabular}
\sourcetable{Notebook 04, d'après le HCP \cite{hcpIndicateurs,hcpENNVM}. Référence : septembre 2022. Les points rapportent la variation des postes couverts au budget initial complet.}
\end{table}


\section{Choc uniforme de \texorpdfstring{5\,\%}{5 \%}}
Le coût du sous-panier augmente de 5\,\% dans les 36 couples profil--périmètre. Dans l'illustration de 5\,000 MAD, le supplément sur les postes couverts varie entre 86,50 et 139,75 MAD dans le périmètre prudent, et entre 223,75 et 245,75 MAD dans l'élargi. À hausse uniforme donnée, cet écart reflète la proportion du budget couverte. Une hausse de 5\,\% sur le budget entier de 5\,000 MAD donnerait 250 MAD.

\section{Choc uniforme suivi du choc alimentaire}
Le supplément du scénario combiné varie entre 215,65 et 403,30 MAD dans le périmètre prudent, et entre 352,90 et 509,30 MAD dans l'élargi. Pour l'ensemble national, le coût du sous-panier prudent augmente de 13,47\,\% (+278,85 MAD et +5,58 points du budget initial), contre 8,76\,\% dans l'élargi (+408,35 MAD et +8,17 points). L'ajout du choc alimentaire représente 175,35 MAD dans les deux périmètres de ce profil, car l'alimentation reçoit le même coefficient d'origine. Parmi les profils, les plus fortes charges alimentaires supplémentaires correspondent aux parts alimentaires initiales les plus élevées. Ainsi, C1 rural consacre 50,2\,\% de son budget au groupe 01, ce qui correspond à 263,55 MAD de supplément alimentaire marginal sur 5\,000 MAD.

Les tableaux \ref{tab:chocs-prudent} et \ref{tab:chocs-elargi} présentent les montants et taux par profil. Les montants illustrent les dépenses supplémentaires nécessaires au maintien des quantités couvertes. Les dépenses réellement engagées dépendraient notamment des substitutions possibles, des prix effectivement rencontrés et des ressources du ménage.
Les contributions au budget initial figurent en annexe~\ref{ann:chocs}. \figureprojet[1.10\linewidth]{04_chocs_hypothetiques.png}{Suppléments en MAD pour un budget initial hypothétique de 5\,000 MAD. Les barres distinguent la hausse uniforme et le supplément alimentaire. « Total » désigne l'ensemble du milieu. Source : notebook 04, pondérations ENNVM \cite{hcpENNVM}.}{fig:chocs}
\clearpage
\begin{table}[!htbp]
\centering\small
\caption{Suppléments sur 5\,000 MAD, périmètre prudent}\label{tab:chocs-prudent}
\begin{tabular}{llrrrr}
\hline
Milieu & Classe & Couvert (\%) & Uniforme (MAD) & Combiné (\%) & Combiné (MAD) \\
\hline
National & C1 & 51{,}9 & 129{,}75 & 14{,}31 & 371{,}25 \\
National & C2 & 48{,}8 & 122{,}00 & 14{,}17 & 345{,}65 \\
National & C3 & 46{,}2 & 115{,}50 & 14{,}00 & 323{,}40 \\
National & C4 & 43{,}0 & 107{,}50 & 13{,}79 & 296{,}50 \\
National & C5 & 36{,}1 & 90{,}25 & 12{,}77 & 230{,}43 \\
National & Ensemble & 41{,}4 & 103{,}50 & 13{,}47 & 278{,}85 \\
Urbain & C1 & 45{,}0 & 112{,}50 & 14{,}12 & 317{,}77 \\
Urbain & C2 & 43{,}5 & 108{,}75 & 13{,}91 & 302{,}48 \\
Urbain & C3 & 42{,}3 & 105{,}75 & 13{,}74 & 290{,}55 \\
Urbain & C4 & 40{,}4 & 101{,}00 & 13{,}47 & 272{,}15 \\
Urbain & C5 & 34{,}6 & 86{,}50 & 12{,}47 & 215{,}65 \\
Urbain & Ensemble & 38{,}6 & 96{,}50 & 13{,}19 & 254{,}53 \\
Rural & C1 & 55{,}9 & 139{,}75 & 14{,}43 & 403{,}30 \\
Rural & C2 & 55{,}1 & 137{,}75 & 14{,}32 & 394{,}48 \\
Rural & C3 & 54{,}0 & 135{,}00 & 14{,}33 & 387{,}00 \\
Rural & C4 & 52{,}6 & 131{,}50 & 14{,}26 & 375{,}10 \\
Rural & C5 & 46{,}5 & 116{,}25 & 14{,}17 & 329{,}40 \\
Rural & Ensemble & 51{,}0 & 127{,}50 & 14{,}26 & 363{,}75 \\
\hline
\end{tabular}
\sourcetable{Simulations hypothétiques. Le taux uniforme vaut 5\,\% dans tous les cas.}
\end{table}

\begin{table}[!htbp]
\centering\small
\caption{Suppléments sur 5\,000 MAD, périmètre élargi}\label{tab:chocs-elargi}
\begin{tabular}{llrrrr}
\hline
Milieu & Classe & Couvert (\%) & Uniforme (MAD) & Combiné (\%) & Combiné (MAD) \\
\hline
National & C1 & 98{,}0 & 245{,}00 & 9{,}93 & 486{,}50 \\
National & C2 & 97{,}1 & 242{,}75 & 9{,}61 & 466{,}40 \\
National & C3 & 95{,}9 & 239{,}75 & 9{,}34 & 447{,}65 \\
National & C4 & 94{,}5 & 236{,}25 & 9{,}00 & 425{,}25 \\
National & C5 & 90{,}1 & 225{,}25 & 8{,}11 & 365{,}43 \\
National & Ensemble & 93{,}2 & 233{,}00 & 8{,}76 & 408{,}35 \\
Urbain & C1 & 97{,}7 & 244{,}25 & 9{,}20 & 449{,}52 \\
Urbain & C2 & 96{,}2 & 240{,}50 & 9{,}03 & 434{,}23 \\
Urbain & C3 & 95{,}1 & 237{,}75 & 8{,}89 & 422{,}55 \\
Urbain & C4 & 93{,}4 & 233{,}50 & 8{,}66 & 404{,}65 \\
Urbain & C5 & 89{,}5 & 223{,}75 & 7{,}89 & 352{,}90 \\
Urbain & Ensemble & 92{,}4 & 231{,}00 & 8{,}42 & 389{,}02 \\
Rural & C1 & 98{,}3 & 245{,}75 & 10{,}36 & 509{,}30 \\
Rural & C2 & 97{,}8 & 244{,}50 & 10{,}25 & 501{,}23 \\
Rural & C3 & 97{,}0 & 242{,}50 & 10{,}20 & 494{,}50 \\
Rural & C4 & 95{,}8 & 239{,}50 & 10{,}09 & 483{,}10 \\
Rural & C5 & 92{,}7 & 231{,}75 & 9{,}60 & 444{,}90 \\
Rural & Ensemble & 95{,}2 & 238{,}00 & 9{,}96 & 474{,}25 \\
\hline
\end{tabular}
\sourcetable{Simulations hypothétiques. Le taux uniforme vaut 5\,\% dans tous les cas.}
\end{table}


\clearpage
\section{Pouvoir d'achat du sous-panier et ressources nominales hypothétiques}
\label{sec:pouvoir-achat}
\subsection{Hypothèses, mesure et seuil de maintien}

Les quatre hypothèses de croissance des ressources nominales sont 0\,\%, +3\,\%, +5\,\% et +10\,\%. Elles s'appliquent au montant effectivement affecté au sous-panier dans l'expérience de simulation. Elles ne proviennent d'aucune série de revenus observés par le HCP. Pour les comparaisons aux prix historiques, chaque hypothèse désigne une variation \emph{totale} entre septembre 2022 et la date considérée. Ces taux ne sont ni annualisés ni composés mois par mois. Pour les deux chocs hypothétiques, ils s'appliquent une seule fois entre l'avant et l'après-choc.

Soient $R_0>0$ les ressources nominales initialement affectées au sous-panier et $C_0>0$ son coût initial. Le nombre d'unités accessibles passe de $R_0/C_0$ à $R_0(1+g_R)/[C_0(1+\pi_h)]$. Pour des taux exprimés en fractions, avec $1+\pi_h>0$, sa variation relative exacte vaut
\begin{equation}
 \Delta_{\mathrm{PA}}
 = \frac{1+g_R}{1+\pi_h}-1
 = \frac{g_R-\pi_h}{1+\pi_h}.
 \label{eq:pouvoir-achat}
\end{equation}
Le résultat est multiplié par 100 dans les tableaux. Le seuil de maintien exact est $g_R^\star=\pi_h$. En dessous, les ressources achètent moins d'unités du sous-panier, et au-dessus elles en achètent davantage. Si le coût diminue, le seuil est négatif. La couverture est reportée pour situer le sous-panier dans le budget total, sans intervenir comme multiplicateur du taux de pouvoir d'achat.

\subsection{Comparaison d'octobre 2026 à septembre 2022}

Pour l'ensemble national, une hausse nominale de 5\,\% donne une variation de $-3{,}86$\,\% dans le périmètre prudent et de $-2{,}64$\,\% dans l'élargi. Le maintien exigerait des hausses nominales de 9,22\,\% et 7,85\,\%, respectivement. Au même taux de ressources, les ensembles urbain et rural affichent respectivement $-3{,}82$\,\% et $-4{,}01$\,\% dans le prudent, puis $-2{,}57$\,\% et $-2{,}93$\,\% dans l'élargi.

Dans chaque milieu, C1 subit une contraction plus forte que C5 pour cette comparaison. Dans le périmètre prudent, le contraste atteint par exemple $-4{,}02$\,\% contre $-3{,}64$\,\% en milieu urbain avec +5\,\% de ressources. Les quintiles sont définis sur les dépenses par personne avec des seuils propres à chaque milieu. Les tableaux \ref{tab:pa-h-prudent} et \ref{tab:pa-h-elargi} détaillent les quatre hypothèses. En octobre 2026, +5\,\% reste inférieur au seuil pour les 18 profils dans les deux périmètres. +10\,\% dépasse tous les seuils de cette comparaison historique.
\figureprojet{05_pouvoir_achat_historique.png}{Variation du nombre d'unités accessibles des sous-paniers en octobre 2026, en référence à septembre 2022. Coûts calculés à partir des prix historiques et croissance nominale hypothétique. Source : notebook 05.}{fig:pa-historique}
\clearpage
\begin{table}[!htbp]
\centering\small
\caption{Octobre 2026 par rapport à septembre 2022, périmètre prudent}\label{tab:pa-h-prudent}
\begin{tabular}{llrrrrrr}
\hline
Milieu & Classe & Couvert & Seuil & 0\,\% & +3\,\% & +5\,\% & +10\,\% \\
\hline
National & C1 & 51{,}9 & 9{,}42 & -8{,}61 & -5{,}87 & -4{,}04 & 0{,}53 \\
National & C2 & 48{,}8 & 9{,}38 & -8{,}58 & -5{,}83 & -4{,}01 & 0{,}57 \\
National & C3 & 46{,}2 & 9{,}36 & -8{,}56 & -5{,}81 & -3{,}99 & 0{,}59 \\
National & C4 & 43{,}0 & 9{,}31 & -8{,}52 & -5{,}77 & -3{,}94 & 0{,}63 \\
National & C5 & 36{,}1 & 9{,}04 & -8{,}29 & -5{,}54 & -3{,}71 & 0{,}88 \\
National & Ensemble & 41{,}4 & 9{,}22 & -8{,}44 & -5{,}70 & -3{,}86 & 0{,}71 \\
\noalign{\vskip 2pt}
Urbain & C1 & 45{,}0 & 9{,}39 & -8{,}59 & -5{,}84 & -4{,}01 & 0{,}56 \\
Urbain & C2 & 43{,}5 & 9{,}35 & -8{,}55 & -5{,}80 & -3{,}97 & 0{,}60 \\
Urbain & C3 & 42{,}3 & 9{,}31 & -8{,}52 & -5{,}77 & -3{,}94 & 0{,}63 \\
Urbain & C4 & 40{,}4 & 9{,}26 & -8{,}48 & -5{,}73 & -3{,}90 & 0{,}67 \\
Urbain & C5 & 34{,}6 & 8{,}97 & -8{,}23 & -5{,}48 & -3{,}64 & 0{,}95 \\
Urbain & Ensemble & 38{,}6 & 9{,}17 & -8{,}40 & -5{,}65 & -3{,}82 & 0{,}76 \\
\noalign{\vskip 2pt}
Rural & C1 & 55{,}9 & 9{,}45 & -8{,}63 & -5{,}89 & -4{,}06 & 0{,}50 \\
Rural & C2 & 55{,}1 & 9{,}41 & -8{,}60 & -5{,}86 & -4{,}03 & 0{,}54 \\
Rural & C3 & 54{,}0 & 9{,}41 & -8{,}60 & -5{,}86 & -4{,}03 & 0{,}54 \\
Rural & C4 & 52{,}6 & 9{,}39 & -8{,}58 & -5{,}84 & -4{,}01 & 0{,}56 \\
Rural & C5 & 46{,}5 & 9{,}34 & -8{,}54 & -5{,}80 & -3{,}97 & 0{,}60 \\
Rural & Ensemble & 51{,}0 & 9{,}38 & -8{,}58 & -5{,}83 & -4{,}01 & 0{,}57 \\
\hline
\end{tabular}
\sourcetable{Notebook 05. Couvert : part du budget initial en pourcentage. Seuil : croissance nominale de maintien en pourcentage. Les quatre dernières colonnes donnent la variation des unités accessibles (en \%).}
\end{table}

\begin{table}[!htbp]
\centering\small
\renewcommand{\arraystretch}{1.10}
\caption{Octobre 2026 par rapport à septembre 2022, périmètre élargi}\label{tab:pa-h-elargi}
\begin{tabular}{llrrrrrr}
\hline
Milieu & Classe & Couverture & Seuil & 0\,\% & +3\,\% & +5\,\% & +10\,\% \\
\hline
National & C1 & 98{,}0 & 8{,}21 & -7{,}59 & -4{,}82 & -2{,}97 & 1{,}65 \\
National & C2 & 97{,}1 & 8{,}16 & -7{,}54 & -4{,}77 & -2{,}92 & 1{,}70 \\
National & C3 & 95{,}9 & 8{,}05 & -7{,}45 & -4{,}68 & -2{,}82 & 1{,}80 \\
National & C4 & 94{,}5 & 7{,}90 & -7{,}32 & -4{,}54 & -2{,}68 & 1{,}95 \\
National & C5 & 90{,}1 & 7{,}65 & -7{,}11 & -4{,}32 & -2{,}46 & 2{,}18 \\
National & \textbf{Ensemble} & 93{,}2 & 7{,}85 & -7{,}27 & -4{,}49 & -2{,}64 & 2{,}00 \\
\noalign{\vskip 3pt}
Urbain & C1 & 97{,}7 & 8{,}07 & -7{,}47 & -4{,}69 & -2{,}84 & 1{,}78 \\
Urbain & C2 & 96{,}2 & 7{,}99 & -7{,}40 & -4{,}62 & -2{,}77 & 1{,}86 \\
Urbain & C3 & 95{,}1 & 7{,}93 & -7{,}34 & -4{,}56 & -2{,}71 & 1{,}92 \\
Urbain & C4 & 93{,}4 & 7{,}81 & -7{,}24 & -4{,}46 & -2{,}60 & 2{,}03 \\
Urbain & C5 & 89{,}5 & 7{,}60 & -7{,}06 & -4{,}27 & -2{,}41 & 2{,}23 \\
Urbain & \textbf{Ensemble} & 92{,}4 & 7{,}77 & -7{,}21 & -4{,}43 & -2{,}57 & 2{,}07 \\
\noalign{\vskip 3pt}
Rural & C1 & 98{,}3 & 8{,}34 & -7{,}70 & -4{,}93 & -3{,}09 & 1{,}53 \\
Rural & C2 & 97{,}8 & 8{,}26 & -7{,}63 & -4{,}86 & -3{,}01 & 1{,}61 \\
Rural & C3 & 97{,}0 & 8{,}29 & -7{,}66 & -4{,}89 & -3{,}04 & 1{,}58 \\
Rural & C4 & 95{,}8 & 8{,}22 & -7{,}60 & -4{,}83 & -2{,}98 & 1{,}64 \\
Rural & C5 & 92{,}7 & 8{,}00 & -7{,}40 & -4{,}63 & -2{,}78 & 1{,}85 \\
Rural & \textbf{Ensemble} & 95{,}2 & 8{,}16 & -7{,}55 & -4{,}77 & -2{,}92 & 1{,}70 \\
\hline
\end{tabular}
\sourcetable{Notebook 05, d'après le HCP \cite{hcpIndicateurs,hcpENNVM}. Couverture et seuil de maintien en \%. Les quatre dernières colonnes donnent la variation des unités accessibles (en \%) par hypothèse nominale.}
\end{table}

\clearpage

\subsection{Chocs hypothétiques et comparaison des périmètres}

Pour la hausse uniforme des prix de 5\,\%, le seuil vaut 5\,\% pour chacun des 36 couples profil--périmètre. Les scénarios nominaux 0\,\%, +3\,\%, +5\,\% et +10\,\% donnent respectivement $-4{,}76$\,\%, $-1{,}90$\,\%, 0\,\% et $+4{,}76$\,\% en unités du sous-panier.

Quand la hausse uniforme est suivie du choc alimentaire supplémentaire, le seuil varie de 12,47\,\% à 14,43\,\% dans le prudent et de 7,89\,\% à 10,36\,\% dans l'élargi. Une hausse nominale de 10\,\% suffit pour 14 des 18 profils dans l'élargi. Les quatre exceptions sont les quintiles ruraux C1 à C4. Elle reste inférieure au seuil pour tous les profils du prudent. Le poste alimentaire représente une part plus élevée du sous-panier prudent, ce qui amplifie l'effet relatif du choc. Les taux des deux périmètres portent sur des compositions et des couvertures différentes.

La figure \ref{fig:pa-chocs} compare les deux chocs. L'annexe \ref{ann:pa-chocs} détaille les quatre variations et le seuil pour les 18 profils après le choc combiné.
\figureprojet{05_pouvoir_achat_chocs.png}{Variations du pouvoir d'achat des sous-paniers sous les chocs de prix et les hypothèses de croissance des ressources. Source : notebook 05.}{fig:pa-chocs}
\clearpage

\subsection{Limites d'interprétation}

Les résultats mesurent une capacité conditionnelle à financer un sous-panier. Les profils sont des agrégats de dépenses et leurs quintiles ne renseignent pas sur le classement des revenus. Le modèle ne dispose ni de trajectoires de revenus effectivement observées ni de quantités suivies permettant d'étudier les substitutions. Il conserve les pondérations de l'ENNVM 2022--2023 à toutes les dates et suppose une évolution proportionnelle des ressources affectées aux postes couverts. Les profils ruraux reçoivent les mêmes variations de prix à collecte urbaine que les profils urbains \cite{hcpIPC}. Les écarts simulés par milieu s'interprètent donc sous cette hypothèse de prix communs. Les résultats ne s'étendent pas aux postes exclus ni au pouvoir d'achat de l'ensemble du budget.

